\section{Clocked traces and quantifier-free internalization}
\label{sec:trace}

Once all finite demands are known, every adaptive construction is represented as a
clocked transition system.  This section gives the central object-proof pattern.

\subsection{Trace specification}

Fix primitive-recursive symbols:
\[
  \operatorname{init}(x),\quad
  \operatorname{step}(t,s,x),\quad
  \operatorname{fuel}(x),\quad
  \operatorname{out}(s,x).
\]
Define the state sequence by primitive recursion:
\[
  S(0,x)=\operatorname{init}(x),\qquad
  S(t+1,x)=\operatorname{step}(t,S(t,x),x).
\]
The state is a numerical code for every accumulator, table, flag, and local certificate
needed by the process.  Let \(I(t,s,x)\) be a bounded formula expressing the invariant,
and let \(Q(x,\operatorname{out}(S(F,x),x))\) be the desired terminal contract, where
\(F=\operatorname{fuel}(x)\).

Booleanize \(I\) and abbreviate
\[
  \widehat I(t,x)\;:\equiv\;\chi_I(t,S(t,x),x)=1.
\]
This is a quantifier-free equation.

\begin{theorem}[QF trace internalization]
\label{thm:trace-internalization}
Suppose \(\iPRA\) derives:
\begin{align*}
 &\widehat I(0,x),\\
 &\widehat I(t,x)\longrightarrow\widehat I(t+1,x),\\
 &\widehat I(F,x)\longrightarrow
   Q(x,\operatorname{out}(S(F,x),x)).
\end{align*}
Then \(\iPRA\) derives
\[
  Q(x,\operatorname{out}(S(F,x),x)).
\]
The construction is uniform in the fixed primitive-recursive syntax and yields an
embedded \(\Prf\,\PRAAx\) term from embedded base, step, and finish derivations.
\end{theorem}

\begin{proof}
Apply quantifier-free induction to \(\widehat I(t,x)\), using the first two derivations.
Universal introduction closes the induction variable, and universal elimination
specializes the result at the term \(F\).  Modus ponens with the terminal derivation
yields the output contract.  At the level of the embedded calculus, these are exactly
\code{PRA.qfInd}, \code{Prf.allE}, \code{Prf.impI}, and \code{Prf.impE}.  No semantic
soundness theorem is needed to construct the proof term.
\end{proof}

The structure \code{QFLoop} and theorems \code{close}, \code{at}, and \code{finish} in
\code{lean/FRNF/Kernel.lean} are a direct source-level implementation of this proof.

\subsection{What belongs in the state?}

A robust state code contains enough information that the step invariant is local.  A
typical record includes:

\begin{itemize}
\item current clock and active/terminal flag;
\item current numerical approximation or exact accumulator;
\item a finite input table or indices into a shared table;
\item a running maximum of requested proof rows;
\item a rank or remaining fuel;
\item local residual, error numerator, or divisibility certificate;
\item finite candidate masks for relational stages; and
\item a checksum or defining-equation witness when useful for efficient reflection.
\end{itemize}

Including more fields can make the invariant simpler.  Since all records are encoded by
pairing and projections, this is a proof-engineering tradeoff rather than a change of
logical strength.

\subsection{Frozen terminal states}

Every clocked process should be made total by freezing after termination:
\[
  \operatorname{step}(t,s,x)=s
  \quad\text{when}\quad\operatorname{terminal}(s,x)=1.
\]
This eliminates partiality and allows a fixed clock larger than the actual stopping
time.  The invariant includes ``terminal states remain terminal and preserve their
output.''  Majorized descent in \cref{sec:descent} supplies a common source of such
clocks.

\subsection{Local reflection versus global semantics}

The trace theorem requires embedded proofs of the base and step equations.  A host
proof that the transition is correct can guide these derivations but cannot replace
them.  In practice the object step proof is assembled from:

\begin{enumerate}[label=(\roman*),leftmargin=2.4em]
\item defining equations for the primitive-recursive transition;
\item Booleanization lemmas for branch guards;
\item arithmetic library lemmas for the updated fields; and
\item local component contracts already present in the proof schedule.
\end{enumerate}

This decomposition is favorable to automation.  The transition itself is fixed and
finite; the compiler can generate the formula and most congruence steps.  Domain
lemmas remain explicit.

\subsection{Variable-length recursion}

A loop whose length is a numerical input is not a schema of separate finite proofs.
It is one primitive recursion and one induction on the clock.  The final theorem may
therefore be genuinely uniform in a matrix size, polynomial degree, path length, or
search bound.  The restriction is that the transition and invariant must be fixed
primitive-recursive syntax.  This is how variable-degree constructions avoid the
universal-evaluator barrier.

\subsection{Nested finite loops}

Nested loops can be flattened into one state machine by storing both counters, or
formalized modularly as components whose local contracts are composed.  Flattening
uses one induction but may produce a large invariant; modularization creates smaller
proofs but requires a proof schedule.  \(\FRNF\) supports both.  The preferred choice
is to isolate mathematically meaningful subroutines and flatten only tightly coupled
adaptive loops.
